% =====================================================================
%  Отчёт по исследованию: KV Cache Quantization × Reasoning Trace Stability
%  Single-file pdflatex source. To compile (from this directory):
%
%      pdflatex report.tex
%      pdflatex report.tex      % run twice for ToC + xrefs
%
%  Required figures (must be in the same directory as this file):
%      accuracy_bars.png
%      divergence_position.png
%
%  Required packages: T2A fontenc, babel-russian, booktabs, graphicx,
%      hyperref, geometry, amsmath/amssymb, listings, xcolor, caption.
%
%  Tested with TeX Live 2023 / MiKTeX 23.x. If hyperref complains about
%  Russian bookmarks, run with:  pdflatex -shell-escape report.tex
% =====================================================================

\documentclass[11pt, a4paper]{article}

\usepackage[T2A]{fontenc}
\usepackage[utf8]{inputenc}
\usepackage[russian, english]{babel}

\usepackage{geometry}
\geometry{margin=2.3cm}

\usepackage{amsmath, amssymb}
\usepackage{graphicx}
\usepackage{booktabs}
\usepackage{array}
\usepackage{caption}
\usepackage{float}
\usepackage{enumitem}
\usepackage{xcolor}
\usepackage[bookmarks=true, colorlinks=true, linkcolor=black,
            citecolor=black, urlcolor=blue]{hyperref}
\usepackage{listingsutf8}

% ---- Listings setup for code/trace excerpts ----
% listingsutf8 + inputencoding=utf8 lets Cyrillic appear inside lstlisting blocks.
\definecolor{lstbg}{rgb}{0.97, 0.97, 0.97}
\lstset{
    basicstyle    = \ttfamily\footnotesize,
    backgroundcolor = \color{lstbg},
    breaklines    = true,
    breakatwhitespace = false,
    columns       = fullflexible,
    frame         = single,
    framesep      = 4pt,
    framerule     = 0pt,
    aboveskip     = 6pt,
    belowskip     = 6pt,
    showstringspaces = false,
    keepspaces    = true,
    extendedchars = true,
    inputencoding = utf8,
    literate=
      {а}{{\selectfont\char224}}1 {б}{{\selectfont\char225}}1
      {в}{{\selectfont\char226}}1 {г}{{\selectfont\char227}}1
      {д}{{\selectfont\char228}}1 {е}{{\selectfont\char229}}1
      {ё}{{\"e}}1                 {ж}{{\selectfont\char230}}1
      {з}{{\selectfont\char231}}1 {и}{{\selectfont\char232}}1
      {й}{{\selectfont\char233}}1 {к}{{\selectfont\char234}}1
      {л}{{\selectfont\char235}}1 {м}{{\selectfont\char236}}1
      {н}{{\selectfont\char237}}1 {о}{{\selectfont\char238}}1
      {п}{{\selectfont\char239}}1 {р}{{\selectfont\char240}}1
      {с}{{\selectfont\char241}}1 {т}{{\selectfont\char242}}1
      {у}{{\selectfont\char243}}1 {ф}{{\selectfont\char244}}1
      {х}{{\selectfont\char245}}1 {ц}{{\selectfont\char246}}1
      {ч}{{\selectfont\char247}}1 {ш}{{\selectfont\char248}}1
      {щ}{{\selectfont\char249}}1 {ъ}{{\selectfont\char250}}1
      {ы}{{\selectfont\char251}}1 {ь}{{\selectfont\char252}}1
      {э}{{\selectfont\char253}}1 {ю}{{\selectfont\char254}}1
      {я}{{\selectfont\char255}}1
      {А}{{\selectfont\char192}}1 {Б}{{\selectfont\char193}}1
      {В}{{\selectfont\char194}}1 {Г}{{\selectfont\char195}}1
      {Д}{{\selectfont\char196}}1 {Е}{{\selectfont\char197}}1
      {Ё}{{\"E}}1                 {Ж}{{\selectfont\char198}}1
      {З}{{\selectfont\char199}}1 {И}{{\selectfont\char200}}1
      {Й}{{\selectfont\char201}}1 {К}{{\selectfont\char202}}1
      {Л}{{\selectfont\char203}}1 {М}{{\selectfont\char204}}1
      {Н}{{\selectfont\char205}}1 {О}{{\selectfont\char206}}1
      {П}{{\selectfont\char207}}1 {Р}{{\selectfont\char208}}1
      {С}{{\selectfont\char209}}1 {Т}{{\selectfont\char210}}1
      {У}{{\selectfont\char211}}1 {Ф}{{\selectfont\char212}}1
      {Х}{{\selectfont\char213}}1 {Ц}{{\selectfont\char214}}1
      {Ч}{{\selectfont\char215}}1 {Ш}{{\selectfont\char216}}1
      {Щ}{{\selectfont\char217}}1 {Ъ}{{\selectfont\char218}}1
      {Ы}{{\selectfont\char219}}1 {Ь}{{\selectfont\char220}}1
      {Э}{{\selectfont\char221}}1 {Ю}{{\selectfont\char222}}1
      {Я}{{\selectfont\char223}}1
      {—}{{---}}1 {–}{{--}}1 {…}{{\ldots}}1
      {«}{{<<}}1 {»}{{>>}}1 {„}{{,,}}1 {“}{{``}}1 {”}{{''}}1
      {≈}{{$\approx$}}1 {≥}{{$\geq$}}1 {≤}{{$\leq$}}1
      {²}{{$^2$}}1 {×}{{$\times$}}1 {·}{{$\cdot$}}1
      {→}{{$\to$}}1 {←}{{$\leftarrow$}}1
      {₀}{{$_0$}}1 {₁}{{$_1$}}1 {₂}{{$_2$}}1
      {§}{{\S}}1 {¶}{{\P}}1 {°}{{$^{\circ}$}}1
      {±}{{$\pm$}}1 {÷}{{$\div$}}1
      {α}{{$\alpha$}}1 {β}{{$\beta$}}1 {γ}{{$\gamma$}}1
      {δ}{{$\delta$}}1 {ε}{{$\varepsilon$}}1 {θ}{{$\theta$}}1
      {λ}{{$\lambda$}}1 {μ}{{$\mu$}}1 {π}{{$\pi$}}1
      {σ}{{$\sigma$}}1 {τ}{{$\tau$}}1 {φ}{{$\varphi$}}1
      {χ}{{$\chi$}}1 {ψ}{{$\psi$}}1 {ω}{{$\omega$}}1
      {Δ}{{$\Delta$}}1 {Σ}{{$\Sigma$}}1
      {∈}{{$\in$}}1 {∉}{{$\notin$}}1 {∞}{{$\infty$}}1
      {∑}{{$\sum$}}1 {∏}{{$\prod$}}1
}

% ---- Convenience macros ----
\newcommand{\code}[1]{\texttt{#1}}
\newcommand{\fp}{$\mathrm{FP}$}
\newcommand{\chisq}{$\chi^{2}$}

% =====================================================================
%  Title
% =====================================================================
\title{\textbf{Влияние квантизации KV-кэша на устойчивость цепочек рассуждений у компактных reasoning-моделей} \\[0.3em]
       \large Отчёт о проделанной работе}
\author{\textit{Студент:} \ldots \\
        \textit{Научный руководитель:} \ldots}
\date{Апрель 2026}

\begin{document}
\maketitle

% =====================================================================
%  Abstract
% =====================================================================
\begin{abstract}
\noindent
Проведено систематическое исследование влияния квантизации KV-кэша
(хранилища ключей и значений механизма внимания во время инференса)
на качество цепочек рассуждений у трёх компактных моделей
(\code{DeepSeek-R1-Distill-Qwen-1.5B/7B}, \code{Qwen3-1.7B}) на $80$
задачах из AIME-24 и MATH-500. Сравнивались пять схем хранения:
\code{bfloat16} (baseline), FP8 (E4M3, E5M2) и HQQ (\,4-бит, 2-бит\,).
Помимо traditional accuracy измерялась \emph{точка первого расхождения}
(First Divergence Point, FDP) между квантизованной и baseline-трассами,
с автоматической классификацией сбоя в этой точке шестикатегорийной
таксономией (Anthropic Claude как судья). Главный результат:
\textbf{семейство Qwen3 устойчиво к FP8 KV-квантизации} (accuracy
$53.7\,\%/51.2\,\%$ против $52.5\,\%$ baseline; \chisq{} между двумя
FP8-схемами $p=0.348$, гипотеза независимости \emph{не отвергается}),
тогда как \textbf{семейство DeepSeek-R1-Distill полностью коллапсирует}
($35{-}60\,\% \to 0\,\%$ под всеми протестированными квантизациями) из-за
доминирования category F (Repetition/loop). Зафиксированы пять случаев
\code{quant\_only} — задач, где квантизация \emph{улучшает} ответ —
все на Qwen3-1.7B, что делает FP8 на этой модели \emph{статистически
взаимозаменяемым} с bf16.
\end{abstract}

\tableofcontents
\bigskip

% =====================================================================
\section{Постановка задачи}
% =====================================================================

\subsection{Контекст и мотивация}

Современные small reasoning models ($1{-}8\,$B параметров) генерируют
chain-of-thought длиной в тысячи токенов. KV-кэш — структура,
содержащая для каждого токена и каждой attention-головы пару тензоров
ключа и значения — занимает основной объём VRAM при инференсе на
длинных контекстах. Например, для \code{DeepSeek-R1-Distill-Qwen-7B}
при генерации $16\,$k токенов KV-кэш занимает $\approx 8{-}12\,$GB.

KV-квантизация (FP8, INT4, INT2) — стандартный приём масштабирования
длинных reasoning-контекстов под фиксированный VRAM-бюджет. Однако в
литературе её влияние оценивается \emph{по итоговой accuracy}\,,
агрегированной по бенчмарку. Этого недостаточно для понимания, \emph{как
именно} квантизация портит рассуждение: меняется ли стратегия решения,
ломается арифметика, или модель срывается в петлю?

\subsection{Гипотезы}

\begin{description}[leftmargin=2.6em, style=nextline]
\item[\textbf{H1 (количественная)}] При понижении точности KV-кэша
    accuracy деградирует монотонно:
    $\mathrm{bf16} \geq \mathrm{FP8} \geq \mathrm{HQQ\text{-}INT4} \geq
    \mathrm{HQQ\text{-}INT2}$.
\item[\textbf{H2 (качественная)}] Тип сбоя (failure mode) зависит от
    агрессивности квантизации: FP8 даёт «арифметические» ошибки,
    HQQ INT2 — «галлюцинации» / срывы стратегии.
\item[\textbf{H3 (архитектурная)}] Семейства моделей различаются
    по чувствительности к квантизации KV; влияние не сводится к
    размеру модели.
\end{description}

\subsection{Цели работы}

\begin{enumerate}[leftmargin=*]
\item Разработать инструмент локализации точки расхождения трасс
(Phase~2) и автоматической классификации сбоев (Phase~3).
\item Получить статистически защищаемые выводы и составить
методологические ограничения для дальнейших работ.
\end{enumerate}

% =====================================================================
\section{Методология}
% =====================================================================

\subsection{Модели и квантизации}

\begin{table}[H]
\centering
\small
\begin{tabular}{lll}
\toprule
\textbf{Модель} & \textbf{Архитектура} & \textbf{HF-идентификатор} \\
\midrule
DS-R1-Distill-Qwen-1.5B & Qwen2 + reasoning distillation & \code{deepseek-ai/DeepSeek-R1-Distill-Qwen-1.5B} \\
Qwen3-1.7B              & Qwen3 (новая архитектура)       & \code{Qwen/Qwen3-1.7B} \\
DS-R1-Distill-Qwen-7B   & Qwen2 + reasoning distillation & \code{deepseek-ai/DeepSeek-R1-Distill-Qwen-7B} \\
\bottomrule
\end{tabular}
\caption{Тестируемые модели.}
\label{tab:models}
\end{table}

\noindent Пять конфигураций KV-кэша:
\begin{itemize}[leftmargin=*]
\item \code{bf16} — baseline, полная точность ($16$ бит);
\item \code{fp8\_e4m3} — FP8, $4$ бита экспоненты, $3$ бита мантиссы (узкий диапазон, высокая точность);
\item \code{fp8\_e5m2} — FP8, $5/2$ (широкий диапазон, ниже точность);
\item \code{hqq\_int4} — HQQ, $4$-битная квантизация ключей и значений;
\item \code{hqq\_int2} — HQQ, $2$-битная квантизация (наиболее агрессивная).
\end{itemize}
FP8 и bf16 запускаются через \textbf{vLLM 0.7.3}; HQQ — через
HuggingFace \code{transformers} с \code{QuantizedCacheConfig(backend="HQQ")},
поскольку vLLM не реализует HQQ.

\subsection{Датасет}

$80$ задач, фиксированный seed: $30$ из AIME-24 (\code{HuggingFaceH4/aime\_2024})
и $50$ из MATH-500 (\code{HuggingFaceH4/MATH-500}). Оба датасета представлены в
обучающих данных всех трёх моделей, что \emph{намеренно}: цель — не измерить
generalisation, а измерить, как квантизация меняет уже выученное поведение
модели в условиях контролируемой нагрузки.
Sampling: \code{temperature=0} (greedy), \code{top\_p=1},
\code{repetition\_penalty=1.05}, \code{max\_tokens=16384}. Greedy
decoding нужен, чтобы любое расхождение трасс было однозначно
отнесено к шуму KV-кэша, а не к семплированию.

\subsection{Конвейер: четыре фазы}

\begin{lstlisting}
Phase 1 (generate)  -> outputs/traces/<model>_<quant>.jsonl
Phase 2 (find_fdp)  -> outputs/fdps/<model>_<quant>.jsonl
Phase 3 (judge)     -> outputs/judgments/<model>_<quant>.jsonl
Phase 4 (analyze)   -> outputs/report.md, outputs/plots/
Phase 5 (post-hoc)  -> outputs/paper/* (этот отчёт)
\end{lstlisting}

\paragraph{Phase 1 --- Generation.} Для каждой пары (model, quant)
генерируются $80$ трасс. Сохраняются: полный текст ответа, токены,
\code{boxed\_answer}, ground truth, \code{finish\_reason}, метаданные.
См. \code{scripts/01\_generate\_traces.py}.

\paragraph{Phase 2 --- FDP detection.} Для каждой задачи берутся две
трассы (bf16 и quantized) одной модели; находится \emph{первая позиция
токена}, где они расходятся (\code{fdp\_token\_idx}). «Косметические»
расхождения (пунктуация, эквивалентные форматирования формул)
отбрасываются с помощью семантической re-sync эвристики на эмбеддингах
\code{sentence-transformers}. Также вычисляется
\code{boxed\_match} $\in$
\{\code{both\_correct}, \code{both\_wrong}, \code{baseline\_only},
\code{quant\_only}, \code{no\_boxed}\}.

\paragraph{Phase 3 --- Judging.} Каждое FDP передаётся Anthropic Claude
(\code{claude-sonnet-4-5}, prompt caching на блоке таксономии).
Возвращаемый JSON содержит категорию A--F, \code{confidence} $\in [0,1]$,
\code{rationale}, \code{affected\_span}. Шесть категорий таксономии:

\begin{itemize}[leftmargin=*, itemsep=0pt]
\item \textbf{A} --- Arithmetic (стратегия верна, но конкретный счёт ошибочен);
\item \textbf{B} --- Logical (посылки верны, вывод некорректен);
\item \textbf{C} --- Strategy-switch (немотивированное переключение метода);
\item \textbf{D} --- Hallucination (изобретение несуществующего факта/теоремы);
\item \textbf{E} --- Premature-termination (трасса оборвалась без ответа);
\item \textbf{F} --- Repetition/loop (один и тот же шаг повторяется $\geq 3$ раз).
\end{itemize}

\paragraph{Phase 4 --- Aggregation.} Глобальный \chisq{} по контингентной
таблице quant $\times$ category и нормализованные failure signatures
(см. \code{outputs/report.md}).

\paragraph{Phase 5 --- Post-hoc paper analysis.} Скрипт
\code{scripts/05\_paper\_analysis.py} читает все артефакты Phase~1\,--\,3
и генерирует восемь дополнительных отчётов в \code{outputs/paper/},
описанных в разделе~\ref{sec:results}.

\subsection{Инфраструктура}

Все вычисления — на одной spot-инстансе \textbf{Vast.ai L40S} (Ada Lovelace,
\code{sm\_89}, $48\,$GB VRAM, $\approx \$1.10{-}1.40$/час). Управление
пайплайном через \code{scripts/run\_all.sh}, диагностика —
\code{scripts/verify\_run.py}. Артефакты выгружены на HuggingFace Hub:
\href{https://huggingface.co/datasets/mikemikemike1111/kv-trace-study}
{\code{mikemikemike1111/kv-trace-study}} ($34$ файла, $79\,$MB).

% =====================================================================
\section{Ход работы и технические препятствия}
% =====================================================================

В этом разделе описаны проблемы, возникшие в процессе сбора данных, и
принятые инженерные решения. Это нужно для оценки достоверности
полученной выборки.

\subsection{Регрессия в \code{transformers} ломает HQQ KV-cache}
\label{sec:transformers-bug}

\paragraph{Симптом.} При запуске \code{qwen3-1.7b} $\times$ \code{hqq\_int4}
на стандартной конфигурации (\code{transformers 4.55.4}, eager attention)
воспроизводится ошибка:

\begin{lstlisting}
RuntimeError: The size of tensor a (N) must match the size of tensor b (2)
              at non-singleton dimension 3
\end{lstlisting}

Стек уходит в \code{eager\_attention\_forward}, шаг
\code{attn\_weights + causal\_mask}. Изначально предположили проблему в
архитектуре Qwen3, но \emph{минимальный воспроизводитель на $15$-токенном
промпте показал ту же ошибку и для DeepSeek-R1-Distill (Qwen2)}.
Следовательно, проблема не специфична для Qwen3.

\paragraph{Диагностика.} Понизили \code{transformers} до \code{4.51.3} —
тот же код работает без изменений. Регрессия в окне \code{4.52} \,$\to$\,
\code{4.55}; точную ревизию не бисектили из-за бюджета времени.

\paragraph{Решение.} Жёсткий пин в \code{requirements.txt}:
\begin{lstlisting}
transformers>=4.51.0,<4.52
\end{lstlisting}
(\code{commits 47b8d3a, 8594ffb}). CI обновлён под тот же пин.

\subsection{SDPA-attention несовместим с HQQ-dequantized тензорами}

После понижения \code{transformers} столкнулись с другой ошибкой в
\code{scaled\_dot\_product\_attention}: KV-тензоры после de-quantization
из HQQ не имеют contiguous strides, нужного Flash-/Memory-Efficient SDPA
ядрам. Решение: принудительно \code{attn\_implementation="eager"} при
загрузке модели (\code{src/kvtrace/generators/hf\_gen.py}).

\subsection{Eager attention слишком медленный для HQQ $\times$ 7B}

После решения 3.1\,--\,3.2 на L40S замерили: один chunk из $4$ задач для
\code{qwen3-1.7b} $\times$ \code{hqq\_int4} занял \textbf{1975 секунд}
($33$ токена/с aggregate, против $615$ токенов/с для bf16 vLLM). Полный
прогон конфигурации ($80$ задач) — $\approx 11$ часов; для 7B-моделей
ещё дольше. Принято \textbf{частичное покрытие HQQ}: оставили
\code{hqq\_int4} и \code{hqq\_int2} только для DeepSeek-1.5B
(уже собранные данные), пропустили Qwen3+HQQ и 7B+HQQ.

\subsection{Разные представления \code{finish\_reason} в vLLM и HF}
\label{sec:hf-finish-reason}

При post-hoc анализе обнаружено, что HF-генератор всегда пишет
\code{finish\_reason="stop"}: \code{model.generate()} не возвращает
per-sample reason. vLLM же различает \code{stop} и \code{length}
корректно. Следствие: для HQQ-ячеек таблица \code{finish\_reason.md}
показывает $100\,\%$ \code{stop} \emph{по построению}. Прокси-метрика —
\code{avg\_gen\_tokens} $\approx \code{max\_tokens}$: если так, сэмпл
скорее всего упёрся в length cap. Полное обсуждение в
разделе~\ref{sec:limitations}, §\,2.3.

\subsection{Бюджет и таймлайн}

\begin{table}[H]
\centering
\small
\begin{tabular}{lrrl}
\toprule
\textbf{Этап} & \textbf{Wall-clock} & \textbf{Стоимость} & \textbf{Заметки} \\
\midrule
Phase 1 (vLLM, bf16+fp8)        & $\sim 3.5\,$ч & $\approx \$4.50$ & 9 конфигураций \\
Phase 1 (HF, HQQ DS-1.5B)       & $\sim 1.0\,$ч & $\approx \$1.30$ & 2 конфигурации \\
Phase 2 (FDP)                   & $\sim 5\,$мин & --- & CPU-bound \\
Phase 3 (judge)                 & $\sim 10\,$мин & $\approx \$3.20$ & $637$ Anthropic-вызовов \\
Phase 4 (aggregate)             & $< 1\,$мин & --- & --- \\
Диагностика 4.55 баг            & $\sim 1.5\,$ч & $\approx \$2.00$ & 3 неудачных запуска + repro \\
Phase 5 (post-hoc)              & --- & --- & local CPU \\
\midrule
\textbf{Итого}                  & $\sim 6\,$ч GPU & $\approx \$11$ & в пределах бюджета \\
\bottomrule
\end{tabular}
\caption{Таймлайн и бюджет.}
\label{tab:budget}
\end{table}

% =====================================================================
\section{Результаты}
\label{sec:results}
% =====================================================================

Все результаты воспроизводимы запуском
\code{python scripts/05\_paper\_analysis.py} после
\code{python scripts/04\_analyze.py}. Каждый пункт ниже —
ссылка на конкретный артефакт.

\subsection{Headline: accuracy по семействам моделей}

\begin{table}[H]
\centering
\small
\begin{tabular}{lrrrrr}
\toprule
\textbf{Модель} & \textbf{bf16} & \code{fp8\_e4m3} & \code{fp8\_e5m2} & \code{hqq\_int4} & \code{hqq\_int2} \\
\midrule
DS-R1-Distill-Qwen-1.5B & $35.0\,\%$ & $0.0\,\%$         & $0.0\,\%$ & $0.0\,\%$ & $0.0\,\%$ \\
Qwen3-1.7B              & $52.5\,\%$ & $\mathbf{53.7\,\%}$ & $51.2\,\%$ & --- & --- \\
DS-R1-Distill-Qwen-7B   & $60.0\,\%$ & $0.0\,\%$         & $0.0\,\%$ & --- & --- \\
\bottomrule
\end{tabular}
\caption{Accuracy на $80$ задачах AIME-24 + MATH-500. «---» = ячейка
не собрана из-за регрессии в \code{transformers} (см. §\,3.1) или
бюджета времени (§\,3.3).}
\label{tab:accuracy}
\end{table}

\begin{figure}[H]
\centering
\includegraphics[width=0.95\textwidth]{accuracy_bars.png}
\caption{Accuracy на AIME-24 + MATH-500 по семействам моделей и
схемам KV-квантизации. Подписи — accuracy и $\Delta$ относительно
bf16. Семейство DeepSeek коллапсирует на любой квантизации; Qwen3
устойчиво.}
\label{fig:accuracy-bars}
\end{figure}

\paragraph{Главный наблюдаемый эффект.} Семейство DeepSeek-R1-Distill
демонстрирует \textbf{полный коллапс} под любой квантизацией KV
(включая самую мягкую FP8-E4M3 на 7B-модели: $60\,\% \to 0\,\%$).
Семейство Qwen3 устойчиво: FP8-E4M3 даже формально \emph{выше} bf16
($53.7$ vs $52.5\,\%$), хотя разница в пределах одной задачи.

\paragraph{Гипотеза H1 (монотонная деградация) отвергнута} в обе
стороны: не наблюдается у Qwen3 (где ничего не деградирует) и нарушается
на DeepSeek-1.5B, где \code{hqq\_int2} (самая агрессивная) даёт ту же
accuracy $0\,\%$, что и \code{fp8\_e4m3} (самая мягкая) — все они ниже
floor возможностей модели.

\paragraph{Гипотеза H3 (архитектурная) подтверждена} наиболее ярко: при
сравнимых размерах ($1.5\,$B vs $1.7\,$B) Qwen3 переносит FP8 без
потерь, а DeepSeek-R1-Distill — нет.

\subsection{Per-model \chisq{} таксономии сбоев (Idea 1)}

Глобальный \chisq{} (\code{outputs/report.md}) даёт $\chi^{2}=66.58$,
$p=1.8\times 10^{-8}$, Cramér's $V=0.187$. Per-model тесты показывают,
что эффект \textbf{не однороден}:

\begin{table}[H]
\centering
\small
\begin{tabular}{lrrrrl}
\toprule
\textbf{Модель} & \chisq{} & \textbf{dof} & \textbf{$p$-value} & \textbf{Cramér's $V$} & \textbf{Интерпретация} \\
\midrule
DS-R1-Distill-Qwen-1.5B & $64.33$ & $9$ & $\mathbf{1.94\times 10^{-10}}$ & $0.259$ & строгая зависимость \\
Qwen3-1.7B              & $5.60$  & $5$ & $\mathbf{0.348}$  & $0.189$ & $H_{0}$ \emph{не отвергается} \\
DS-R1-Distill-Qwen-7B   & $9.54$  & $2$ & $8.50\times 10^{-3}$ & $0.244$ & e5m2 $\to$ $100\,\%$ F \\
\bottomrule
\end{tabular}
\caption{Per-model \chisq{} тесты независимости failure-mode распределения от
схемы квантизации. Для Qwen3 две FP8-схемы статистически взаимозаменяемы.}
\label{tab:per-model-chi2}
\end{table}

\paragraph{Содержательная интерпретация.}
\begin{itemize}[leftmargin=*]
\item \textbf{DeepSeek-1.5B}: разные quant дают разные паттерны сбоев.
Standardized residual в категории D (Hallucination) для
\code{fp8\_e4m3} составляет $+4.09$ (сильный outlier) — \code{fp8\_e4m3}
толкает модель в галлюцинации, тогда как \code{fp8\_e5m2} толкает в
петли (F).
\item \textbf{Qwen3-1.7B}: две FP8-схемы статистически взаимозаменяемы.
Это важный практический вывод: при выборе между E4M3 и E5M2 на Qwen3
разница в пределах шума.
\item \textbf{DeepSeek-7B}: \code{e5m2} даёт $100\,\%$ Repetition/loop
($80/80$); \code{e4m3} даёт $89\,\%$ F + $9\,\%$ D + $2\,\%$ E. Разница в
\emph{типе} сбоя при одинаковом коллапсе accuracy.
\end{itemize}

\subsection{Аномалия \code{quant\_only} на Qwen3 (Idea 3)}

\code{boxed\_match = "quant\_only"} — задачи, где quantized-ветвь дала
правильный ответ, а bf16-baseline — нет. Это нарушает наивное ожидание
«квантизация может только ухудшать».

Найдено \textbf{$5$ таких случаев}, все на Qwen3-1.7B:
\begin{itemize}[itemsep=0pt]
\item $2$ на \code{qwen3-1.7b $\times$ fp8\_e4m3} (problems 8, 14 из AIME-24);
\item $3$ на \code{qwen3-1.7b $\times$ fp8\_e5m2}.
\end{itemize}

Один представительный кейс приведён в Приложении~\ref{app:case1}; полный
side-by-side анализ всех пяти случаев — в \code{outputs/paper/quant\_only\_deepdives.md}.
Типичный паттерн: bf16-модель \emph{застревает} в избыточной верификации
найденного решения и не успевает выписать \code{\textbackslash{}boxed\{\}}
до \code{max\_tokens}, тогда как quantized-ветвь после расхождения
\emph{короче} и успевает. Судья классифицирует расхождение как
Strategy-switch (категория C), но конкретно эти switch'и оказались
полезными.

Параллельно зафиксированы и \code{baseline\_only}-случаи (bf16 правильно,
quant — нет): $1$ на e4m3, $4$ на e5m2 для Qwen3. Сумма квалифицирующих
расхождений примерно симметрична, что подтверждает интерпретацию
\textbf{квантизация как стохастическое возмущение} на Qwen3, а не как
систематическое ухудшение.

\subsection{Позиция расхождения в трассе (Extra 5)}

Нормированная позиция $\mathrm{FDP\_idx} / \mathrm{baseline\_length}$:

\begin{table}[H]
\centering
\small
\begin{tabular}{llrrr}
\toprule
\textbf{Модель} & \textbf{Quant} & \textbf{n} & \textbf{median} & \textbf{mean} \\
\midrule
DS-R1-Distill-Qwen-1.5B & \code{fp8\_e4m3} & $80$ & $0.00$ & $0.01$ \\
DS-R1-Distill-Qwen-1.5B & \code{fp8\_e5m2} & $80$ & $0.00$ & $0.01$ \\
DS-R1-Distill-Qwen-7B   & \code{fp8\_e4m3} & $80$ & $0.00$ & $0.05$ \\
DS-R1-Distill-Qwen-7B   & \code{fp8\_e5m2} & $80$ & $0.00$ & $0.00$ \\
Qwen3-1.7B              & \code{fp8\_e4m3} & $80$ & $\mathbf{0.07}$ & $0.16$ \\
Qwen3-1.7B              & \code{fp8\_e5m2} & $80$ & $0.03$ & $0.13$ \\
\bottomrule
\end{tabular}
\caption{Распределение позиции первого расхождения. Median вблизи $0$ =
расхождение в первых процентах трассы.}
\label{tab:divergence-position}
\end{table}

\begin{figure}[H]
\centering
\includegraphics[width=\textwidth]{divergence_position.png}
\caption{Гистограммы позиции FDP по (model, quant). Вертикальная
пунктирная линия — медиана. DeepSeek-модели (строки 1, 3) показывают
расхождение почти сразу (всё сосредоточено в $[0, 0.05]$); Qwen3
(строка 2) демонстрирует существенный разброс по всей трассе, медиана
$\approx 0.07$. «No data» = ячейки, не собранные из-за §\,3.1 / §\,3.3.}
\label{fig:divergence-position}
\end{figure}

\paragraph{Содержательно (с поправкой на средние длины трасс).}
\begin{itemize}[leftmargin=*, itemsep=0pt]
\item DeepSeek модели расходятся в среднем за \emph{первые $1\,\%$} токенов
— $\approx 170$ токенов на 1.5B и $\approx 300$ токенов на 7B. Шум KV-кэша
попадает в reasoning-цепочку почти сразу.
\item Qwen3-1.7B расходится в среднем на $7\,\% / 13\,\%$ длины трассы —
$\approx 540 / \approx 1000$ токенов. Модель «удерживает» рассуждение
значительно дольше до проявления эффекта квантизации.
\end{itemize}
Это, по-видимому, ключевой механизм наблюдаемой устойчивости Qwen3:
расхождение происходит позже, и ёмкости оставшейся трассы хватает на
самокоррекцию.

\subsection{Token efficiency (Extra 6)}

\begin{table}[H]
\centering
\small
\begin{tabular}{llrrrr}
\toprule
\textbf{Модель} & \textbf{Quant} & \textbf{n} & \textbf{avg\_gen} & \textbf{$\Delta$ vs bf16} & \textbf{\% length} \\
\midrule
DS-R1-Distill-Qwen-1.5B & \code{bf16}      & $80$ & $17077$ & ---     & $48\,\%$ \\
DS-R1-Distill-Qwen-1.5B & \code{fp8\_e4m3} & $80$ & $16180$ & $-5.2\,\%$ & $\mathbf{98\,\%}$ \\
DS-R1-Distill-Qwen-1.5B & \code{fp8\_e5m2} & $80$ & $16094$ & $-5.8\,\%$ & $\mathbf{89\,\%}$ \\
Qwen3-1.7B              & \code{bf16}      & $80$ & $7748$  & ---     & $22\,\%$ \\
Qwen3-1.7B              & \code{fp8\_e4m3} & $80$ & $7949$  & $+2.6\,\%$ & $26\,\%$ \\
Qwen3-1.7B              & \code{fp8\_e5m2} & $80$ & $8099$  & $+4.5\,\%$ & $25\,\%$ \\
DS-R1-Distill-Qwen-7B   & \code{bf16}      & $80$ & $5997$  & ---     & $19\,\%$ \\
DS-R1-Distill-Qwen-7B   & \code{fp8\_e4m3} & $80$ & $15960$ & $\mathbf{+166\,\%}$ & $\mathbf{92\,\%}$ \\
DS-R1-Distill-Qwen-7B   & \code{fp8\_e5m2} & $80$ & $15805$ & $\mathbf{+164\,\%}$ & $\mathbf{92\,\%}$ \\
\bottomrule
\end{tabular}
\caption{Token efficiency. На DeepSeek-7B FP8 не «ускоряет» инференс,
а \emph{в $2.66$ раза удлиняет} среднюю генерацию из-за петель. На
Qwen3 длина почти не меняется — квантизация выгодна по памяти и не
вредна по латентности.}
\label{tab:token-efficiency}
\end{table}

\subsection{Распределение \code{finish\_reason} (Extra 8)}

Прямая верификация механизма категории F: на DeepSeek-7B обе FP8-схемы
дают $92\,\%$ \code{finish\_reason=length} (vLLM-репорт), что напрямую
соответствует $80/80$ категории F (Repetition/loop) в
\code{report.md}. Разница в типе сбоя между моделями — это разница в
том, \emph{где} обрывается генерация: DeepSeek уходит в петлю и
расходует \code{max\_tokens}; Qwen3 нормально терминируется с EOS.

\subsection{Валидация таксономии: confidence судьи (Extra 7)}

\begin{table}[H]
\centering
\small
\begin{tabular}{lrrrrr}
\toprule
\textbf{Категория} & \textbf{n} & \textbf{mean\_conf} & \textbf{std} & \textbf{min} & \textbf{$n_{<0.5}$} \\
\midrule
A (Arithmetic)             & $21$  & $0.80$ & $0.10$ & $0.55$ & $0$ \\
B (Logical)                & $17$  & $0.79$ & $0.08$ & $0.60$ & $0$ \\
C (Strategy-switch)        & $64$  & $0.74$ & $0.10$ & $0.55$ & $0$ \\
D (Hallucination)          & $51$  & $0.84$ & $0.07$ & $0.72$ & $0$ \\
E (Premature-termination)  & $72$  & $0.83$ & $0.13$ & $0.40$ & $1$ \\
F (Repetition/loop)        & $412$ & $\mathbf{0.95}$ & $0.05$ & $0.72$ & $0$ \\
\bottomrule
\end{tabular}
\caption{Средняя confidence судьи (Anthropic Claude) по категориям.
F (петля) распознаётся почти идеально (mean = $0.95$); C
(Strategy-switch) — наименее уверенно (mean = $0.74$).}
\label{tab:judge-confidence}
\end{table}

Все категории имеют $< 5$ случаев confidence $< 0.5$, что валидирует
таксономию как достаточно различимую для целей агрегации.

\subsection{Глобальная статистика расхождений (Extra 9)}

\begin{table}[H]
\centering
\scriptsize
\setlength{\tabcolsep}{4pt}
\begin{tabular}{llrrrrrr}
\toprule
\textbf{Модель} & \textbf{Quant} & \textbf{both\_correct} & \textbf{both\_wrong} & \textbf{quant\_only} & \textbf{baseline\_only} & \textbf{no\_boxed} \\
\midrule
DS-R1-Distill-Qwen-1.5B & \code{fp8\_e4m3} & $0$  & $15$ & $0$         & $28$ & $37$ \\
DS-R1-Distill-Qwen-1.5B & \code{fp8\_e5m2} & $0$  & $14$ & $0$         & $28$ & $38$ \\
DS-R1-Distill-Qwen-1.5B & \code{hqq\_int4} & $0$  & $14$ & $0$         & $28$ & $38$ \\
DS-R1-Distill-Qwen-1.5B & \code{hqq\_int2} & $0$  & $14$ & $0$         & $28$ & $38$ \\
Qwen3-1.7B              & \code{fp8\_e4m3} & $41$ & $22$ & $\mathbf{2}$ & $1$  & $14$ \\
Qwen3-1.7B              & \code{fp8\_e5m2} & $38$ & $23$ & $\mathbf{3}$ & $4$  & $12$ \\
DS-R1-Distill-Qwen-7B   & \code{fp8\_e4m3} & $0$  & $17$ & $0$         & $48$ & $15$ \\
DS-R1-Distill-Qwen-7B   & \code{fp8\_e5m2} & $0$  & $17$ & $0$         & $48$ & $15$ \\
\bottomrule
\end{tabular}
\caption{Распределение \code{boxed\_match} по (model, quant). Все $5$
случаев \code{quant\_only} — на Qwen3-1.7B.}
\label{tab:fdp-rate}
\end{table}

Видно, что для DeepSeek \code{both\_correct = 0} во всех клетках: ни в
одном случае quantized-ветвь не воспроизвела правильный ответ, даже
если baseline его получил. Для Qwen3 \code{both\_correct} доминирует
($38{-}41$ из $80$): квантизация в большинстве задач сохраняет
правильность.

% =====================================================================
\section{Выводы}
% =====================================================================

\begin{enumerate}[leftmargin=*]
\item \textbf{H3 подтверждена.} Семейство модели — главный фактор
чувствительности к квантизации KV-кэша при сравнимом размере.
Архитектурно более новый Qwen3-1.7B полностью переносит FP8;
distilled DeepSeek-R1 коллапсирует уже на E4M3.

\item \textbf{H1 опровергнута.} Деградация не монотонна по
агрессивности квантизации. На DeepSeek-1.5B accuracy $= 0\,\%$ уже на
самой мягкой схеме, и далее не падает (нечему). На Qwen3 accuracy
не падает вовсе.

\item \textbf{H2 частично подтверждена.} Тип сбоя действительно
зависит от quant: на DeepSeek-1.5B \code{fp8\_e4m3} даёт $+4.09$
standardized residual в категории D (галлюцинации), \code{fp8\_e5m2}
уходит в F (петли); \code{hqq\_int2} уходит в E (преждевременная
остановка). Однако сила эффекта (Cramér's $V = 0.26$) —
small-to-moderate, и для воспроизводимого утверждения на уровне
«E4M3 $\to$ арифметика, INT2 $\to$ галлюцинации» нужен сэмпл-сайз
больше $80$ задач.

\item \textbf{Качественное наблюдение.} Позиция расхождения
($\approx 1\,\%$ длины для DeepSeek vs $\approx 10\,\%$ для Qwen3)
— кандидатный механизм для устойчивости Qwen3. Гипотеза для будущей
работы: устойчивые модели «амортизируют» шум KV в первых сотнях
токенов рассуждения.
\end{enumerate}

% =====================================================================
\section{Ограничения}
\label{sec:limitations}
% =====================================================================

Полный список и обоснование — в \code{outputs/paper/limitations.md}.
Кратко:
\begin{itemize}[leftmargin=*]
\item \textbf{Покрытие}: $4$ из $15$ ячеек (model $\times$ quant) не
собраны из-за регрессии в \code{transformers $\geq 4.55.0$} (см. §\,3.1)
и/или wall-clock budget на eager-attention пути (§\,3.3). Все эти $4$ —
HQQ $\times$ \{Qwen3-1.7B, DeepSeek-7B\}.
\item \textbf{Sampling}: только greedy \code{temperature=0}; нет оценки
variance.
\item \textbf{\code{repetition\_penalty=1.05}} для всех конфигураций —
для DeepSeek может быть слишком мягким; вероятно, при $1.10{-}1.15$
категория F уступит место «настоящим» арифметико-логическим ошибкам,
что изменит таблицу таксономии.
\item \textbf{\code{finish\_reason} в HF-генераторе} некорректен (всегда
\code{stop}); требует фикса для будущих HQQ-прогонов (см. §\,3.4).
\item \textbf{Один судья} (Anthropic Claude); $\kappa$ с альтернативным
судьёй не считалась.
\item \textbf{Бенчмарки контаминированы} в обучающих данных моделей —
внешняя валидность ограничена.
\end{itemize}

% =====================================================================
\section{Будущая работа}
% =====================================================================

Конкретные эксперименты, ранжированные по приоритету и стоимости:

\begin{enumerate}[leftmargin=*]
\item \textbf{Sweep по \code{repetition\_penalty}} на DeepSeek $\times$
FP8 ($4$ конфигурации $\times 80$ задач $\times 4$ значения penalty
$= 1280$ trials, $\approx \$8$ GPU). Цель: проверить, проявляются ли
«настоящие» сбои при подавлении петель.

\item \textbf{Дозбор HQQ ячеек} после фикса в \code{transformers}
($4$ ячейки, $\approx \$15$ GPU при возврате к vLLM-friendly stack или
$\approx \$50$ при текущем eager path).

\item \textbf{Multi-seed ($n=8$) на \code{temperature=0.6}} для оценки
variance — $\approx \$30$ GPU.

\item \textbf{Cross-architecture replication}: добавить
\code{Llama-3.2-3B} и \code{Qwen2.5-3B} для проверки, является ли
«Qwen3 robust, DeepSeek collapses» явлением семейства или специфики
distillation.

\item \textbf{Inter-rater agreement}: 100-divergence subset через
GPT-4o, Cohen's $\kappa$.

\item \textbf{Per-layer KV-quant ablation}: квантизовать только первые
$N$ или только последние $N$ слоёв — поиск «уязвимого» региона.
\end{enumerate}

% =====================================================================
\section{Воспроизводимость и артефакты}
% =====================================================================

\subsection{Код}

Репозиторий: https://github.com/neuraCollab/kv-cache-small-distill-errors. Ключевые скрипты:
\begin{itemize}[leftmargin=*, itemsep=0pt]
\item \code{scripts/01\_generate\_traces.py} — Phase 1
\item \code{scripts/02\_find\_fdps.py} — Phase 2
\item \code{scripts/03\_judge\_fdps.py} — Phase 3
\item \code{scripts/04\_analyze.py} — Phase 4
\item \code{scripts/05\_paper\_analysis.py} — Phase 5 (post-hoc)
\item \code{scripts/verify\_run.py} — диагностика после прогона
\item \code{scripts/run\_all.sh} — orchestrator
\end{itemize}

\subsection{Зависимости}

\code{requirements.txt} фиксирует все пины. Критические:
\begin{lstlisting}
vllm        >= 0.7.3, < 0.9
transformers>= 4.51.0, < 4.52   # см. § 3.1
hqq         >= 0.2.5
anthropic   == 0.40.0
\end{lstlisting}

\subsection{Артефакты}

Все $11$ трасс + $8$ FDP + $8$ judgments + \code{report.md} + \code{paper/*}
выгружены на HuggingFace Hub:
\href{https://huggingface.co/datasets/mikemikemike1111/kv-trace-study}
{\code{mikemikemike1111/kv-trace-study}} ($34$ файла, $79\,$MB).

\subsection{CI и тесты}

GitHub Actions workflow \code{.github/workflows/ci.yml} запускает на
каждый PR: \code{ruff}, \code{mypy}, \code{pytest} (CPU-only маркеры) с
покрытием $\geq 85\,\%$. Live-API тесты (Anthropic) и GPU-only тесты
помечены \code{@pytest.mark.live\_api}, \code{@pytest.mark.gpu} и
пропускаются в CI.

% =====================================================================
\appendix
\section{Представительный \code{quant\_only} кейс}
\label{app:case1}
% =====================================================================

\textit{Полный набор из $5$ кейсов с side-by-side трассами —
в \code{outputs/paper/quant\_only\_deepdives.md}. Здесь приведён один,
наиболее показательный.}

\paragraph{Case 1: \code{qwen3-1.7b} $\times$ \code{fp8\_e4m3} $\times$ problem 8 (AIME-24).}
\begin{itemize}[leftmargin=*, itemsep=0pt]
\item \textbf{FDP token idx}: $554$
\item \textbf{Ground truth}: $25$
\item \textbf{Baseline boxed}: \code{None} (не успел дописать)
\item \textbf{Quant boxed}: $25$
\item \textbf{Judge category}: \textbf{C (Strategy-switch)}, confidence $0.85$
\item \textbf{Judge rationale}: «The quantized trace abandons the clean
substitution approach (\,$k = \log_{y}(x)$\,) that was working and switches to
substituting $y = x^{10/x}$ and then yet another approach, without
motivation. The baseline continues directly to the answer using the
established substitution.»
\end{itemize}

\paragraph{Задача.}
\begin{quote}
\itshape
There exist real numbers $x$ and $y$, both greater than $1$, such that
$\log_{x}(y^{x}) = \log_{y}(x^{4y}) = 10$. Find $xy$.
\end{quote}

\paragraph{Общий префикс рассуждения (последняя часть, идентичная в обеих ветвях).}
\begin{lstlisting}
4y * log_y(x) = 10.

Now, I have two equations:

1) x * log_x(y) = 10
2) 4y * log_y(x) = 10

Hmm, maybe I can express log_x(y) in terms of log_y(x)? Because
log_x(y) is the reciprocal of log_y(x), right? ... So, log_x(y) =
1 / log_y(x). ...

So, if I let log_y(x) be some variable, say, let me denote log_y(x)
as k. Then log_x(y) would be 1/k.

So, substituting back into the equations:
First equation:  x * (1/k) = 10  =>  x/k = 10  =>  x = 10k.
Second equation: 4y * k = 10     =>  yk = 10/4 = 5/2.

So now, we have x
\end{lstlisting}

\paragraph{Baseline (bf16) после расхождения.}
\begin{lstlisting}
... So now, we have x = 10k and y = (5/2)/k. Therefore,
xy = 10k * (5/2)/k = 10*(5/2) = 25.

Wait, so is it 25? That seems straightforward, but let me check
if there are any constraints or if I made a mistake.

Wait, but hold on. Let me make sure that my substitution is
correct. Let me retrace.

Let me define k = log_y(x). Then ...
[trace продолжается верификацией; модель не успевает выписать
 \boxed{} до max_tokens = 16384]
\end{lstlisting}

\paragraph{Quantized (\code{fp8\_e4m3}) после расхождения.}
\begin{lstlisting}
... but maybe.

Now moving on to the second equation: log_y(x^{4y}) = 10.
Again, using the same logarithm property:
log_y(x^{4y}) = 4y * log_y(x) = 10.

Similarly, if I let log_y(x) = m, then y^m = x. And from the
equation above, 4y * m = 10. So m = 10/(4y) = 5/(2y).

But also, log_y(x) = m = ln(x)/ln(y). But maybe I can relate this
to the previous expression?

Wait, earlier I had from the first equation that y = x^{10/x}.
... [более короткий путь, успевает дойти до boxed{25}]
\end{lstlisting}

\paragraph{Интерпретация.} Baseline \emph{избыточно верифицирует} уже
найденное решение и расходует токены на retracing своих собственных
выкладок. Quantized-ветвь после расхождения переходит на другую
параметризацию ($y = x^{10/x}$), которая локально кажется судье
немотивированной (категория C), но \emph{короче} ведёт к ответу — и
успевает выписать \code{\textbackslash{}boxed\{25\}} до \code{max\_tokens}.
Случай иллюстрирует, что метрика accuracy в режиме длинных reasoning-трасс
существенно зависит от длины генерации, а квантизация — стохастический
сдвиг этой длины.

% =====================================================================
\section{Список supplementary артефактов}
\label{app:artifacts}
% =====================================================================

Все артефакты в директории \code{outputs/paper/}:

\begin{table}[H]
\centering
\small
\begin{tabular}{ll}
\toprule
\textbf{Файл} & \textbf{Содержание} \\
\midrule
\code{report.tex}                  & Этот документ \\
\code{supervisor\_report.md}       & Markdown-версия для быстрого чтения \\
\code{per\_model\_chi2.md, .csv}   & Per-model \chisq{} и Cramér's $V$ (Idea 1) \\
\code{accuracy\_bars.png, .svg}    & Bar chart accuracy (Idea 2; figure~\ref{fig:accuracy-bars}) \\
\code{quant\_only\_deepdives.md}   & Все $5$ \code{quant\_only} кейсов с трассами (Idea 3) \\
\code{limitations.md}              & Полный список ограничений (Idea 4) \\
\code{divergence\_position.png}, .svg, .md & Гистограммы и таблица позиции FDP (Extra 5; figure~\ref{fig:divergence-position}) \\
\code{token\_efficiency.md}        & Token efficiency (Extra 6; table~\ref{tab:token-efficiency}) \\
\code{judge\_confidence.md}        & Validation таксономии (Extra 7; table~\ref{tab:judge-confidence}) \\
\code{finish\_reason.md}           & Распределение \code{finish\_reason} (Extra 8) \\
\code{fdp\_rate.md}                & Per-model FDP rate (Extra 9; table~\ref{tab:fdp-rate}) \\
\bottomrule
\end{tabular}
\caption{Inventory artefacts.}
\end{table}

\end{document}
